\documentclass[11pt,a4paper]{article}

\usepackage[utf8]{inputenc}
\usepackage[T1]{fontenc}
\usepackage{amsmath,amssymb,amsthm,mathtools}
\usepackage{graphicx}
\graphicspath{{../../validation-incompleteness/figures/}{./}}
\usepackage{booktabs}
\usepackage{array}
\usepackage{enumitem}
\usepackage{float}
\usepackage{siunitx}
\usepackage{microtype}
\usepackage{lmodern}
\usepackage[margin=2.4cm]{geometry}
\usepackage[hidelinks]{hyperref}
\usepackage{cleveref}
\usepackage{natbib}

\theoremstyle{plain}
\newtheorem{theorem}{Theorem}[section]
\newtheorem{lemma}[theorem]{Lemma}
\newtheorem{proposition}[theorem]{Proposition}
\newtheorem{corollary}[theorem]{Corollary}
\theoremstyle{definition}
\newtheorem{definition}[theorem]{Definition}
\newtheorem{assumption}[theorem]{Assumption}
\newtheorem{construction}[theorem]{Construction}
\theoremstyle{remark}
\newtheorem{remark}[theorem]{Remark}
\newtheorem{example}[theorem]{Example}

% ---------------------------------------------------------------------------
% Notation local to this paper.
% ---------------------------------------------------------------------------
\newcommand{\Model}{\mathcal{M}}          % a model (schema + axioms)
\newcommand{\Corpus}{\mathcal{D}}         % a corpus (ground facts)
\newcommand{\Engine}{\mathcal{E}}         % an engine (compiler + reasoner)
\newcommand{\Feat}{\mathcal{F}}           % the feature alphabet
\newcommand{\req}{\varrho}                % feature requirement of a question
% \sup is a LaTeX operator already, so this must be \renewcommand; a
% \newcommand here raises "Command \sup already defined".
\renewcommand{\sup}{\operatorname{Sup}}   % declared capability set
\newcommand{\Ans}{\mathsf{A}}             % answer set
\newcommand{\Verd}{\mathsf{V}}            % verdict
\newcommand{\vans}{\textsc{answered}}
\newcommand{\vexp}{\textsc{cannot-express}}
\newcommand{\vder}{\textsc{not-derivable}}
\newcommand{\vsurf}{\textsc{no-query-surface}}
\newcommand{\vtime}{\textsc{timeout}}
\newcommand{\vvac}{\textsc{control-vacuous}}
\newcommand{\vinert}{\textsc{constraint-inert}}
% Named \lowr because \lower is a TeX box primitive: overriding it would
% break any package that relies on it.
\newcommand{\lowr}{\Lambda}               % lowering map
\newcommand{\qq}{q}
\newcommand{\Ques}{\mathcal{Q}}

\begin{document}

\title{%
\textbf{The Expectation of Incompleteness}\\[0.4em]
\large A Formal Account of Knowledge Graphs as Instruments That Report
What They Cannot Say}

\author{%
Kundai Farai Sachikonye\\[0.3em]
\normalsize Experimental Bioinformatics, TUM School of Life Sciences\\
\normalsize Maximus-von-Imhof-Forum 3, 85354 Freising, Germany\\
\normalsize \texttt{kundai.sachikonye@bitspark.com}}

\date{\today}

\maketitle

\begin{abstract}
A causal knowledge graph over any biochemical corpus will fail to account for
some of the reactions in that corpus. This is not a defect to be engineered
away; it is the expected state of every such artefact, and the interesting
question is not how much is covered but whether the system can say, at each
question it is asked, \emph{which} of its several independent limitations is
responsible for the failure. We argue that a knowledge graph should be built
as a measuring instrument whose primary output is a verdict rather than an
answer set, and we develop the formal apparatus for this from first
principles. We define three logically independent predicates over a
question---\emph{expressible} in a model, \emph{derivable} from a corpus under
that model's semantics, and \emph{answerable} by a given engine---and prove
that no two of them entail the third, so that the customary practice of
reporting a single empty answer set conflates three distinct claims. We
introduce a verdict space in which \vexp, \vder{} and \vsurf{} are outcomes
with the same evidential standing as \vans, and prove a
\emph{non-degeneracy} property: no verdict other than \vans{} may carry an
empty answer set, and \vans{} may not carry one either unless emptiness is
itself certified. We then show that a system's declared capability set is a
claim about its \emph{compiler} and not about its underlying formalism, that
the two claims come apart in both directions, and that only the compiler-level
reading admits a decidable containment check. Because a constraint may be
lowered into a target language in a form that is syntactically present and
semantically inert, we prove that syntactic and semantic vacuity checks are
mutually irredundant---neither subsumes the other---and that a system carrying
only one of them can certify a result produced by a constraint that did no
work. Finally we prove a \emph{single-derivation} property: when every view of
a result is computed from one derivation, two views of the same system cannot
disagree, and we show that duplicated derivation admits contradictory
publication with every schema check passing. Seventeen computational
experiments accompany the paper. Each is stated as a proposition, given a
control able to fail, and reported with its verdict; several are designed so
that the interesting outcome is a refusal. Everything required to read this
paper is developed here from first principles, with no reliance on prior work
by the author.
\end{abstract}

\tableofcontents

\section{Introduction}

\subsection{The failure mode that motivates this paper}

Consider a system that answers questions over a biochemical corpus by
compiling each question into some logical form, running a reasoner, and
returning the resulting set of entities. Ask it: \emph{which reactions have
exactly four participants?} Suppose the corpus contains three such reactions
and the system returns the empty set.

There are at least five distinct situations in which that output is produced.
The model may lack cardinality constructs altogether, so the question was never
representable. The model may have them, but under an open-world reading the
count is not entailed, because nothing rules out a fourth participant that
merely was not recorded. The corpus may genuinely contain no such reaction.
The compiler may have silently dropped the cardinality restriction while
lowering the question, so that a syntactically well-formed query ran and
matched nothing. Or the reasoner may have exhausted a search budget and
returned partial results without saying so.

These five situations warrant five different actions on the part of whoever
receives the answer, and exactly one output is produced for all of them. The
output is not wrong in any way that inspection reveals: it is well formed, it
is of the correct type, it is returned promptly, and in the second case it is
even \emph{correct} under the semantics the system nominally implements. It is
simply uninformative in a way that is invisible.

This paper is about building systems in which those five situations produce
five different outputs, and about what that requirement does to the
architecture. Our thesis is that the requirement is not a reporting nicety
bolted onto a query engine but a structural constraint that determines what
the system is: a knowledge graph built this way is an \emph{instrument}, and
its output at each question is a \emph{verdict}, of which an answer set is one
possible constituent.

\subsection{Why incompleteness is the expected state}

It is worth being explicit that we do not regard the incompleteness of a
biochemical knowledge graph as a temporary condition. Three independent
sources guarantee it.

The first is expressive. Any model is a fragment of some logic, and every
decidable fragment of first-order logic omits constructions that some question
will need~\citep{BaaderCalvanese2017}. The trade-off between expressive power
and the tractability of reasoning is not an engineering artefact but a
theorem-level fact about the fragments~\citep{Donini2003,Grau2008}. A model
chosen so that reasoning terminates is a model in which some questions are not
statable.

The second is empirical. Curated biochemical corpora are assembled by human
effort under finite budget from an open and growing literature. Reaction
databases record what has been curated, not what
exists~\citep{Alcantara2012,Bansal2022}, and chemical ontologies partition
molecules along axes chosen for chemical rather than metabolic
convenience~\citep{Degtyarenko2008,Hastings2016}, so that species a biochemist
would group together may share no informative common ancestor. Absence in a
corpus is not evidence of absence in the world, and treating it as such is the
closed-world error.

The third is computational. Even where a question is both statable and
entailed, the reasoner may not finish. Description-logic reasoning is
intractable in the worst case for the expressive
profiles~\citep{Motik2009Hypertableau,Glimm2014HermiT}, and the profiles
designed for tractability~\citep{Motik2009Profiles,Kazakov2014ELK} achieve it
precisely by dropping constructs. A question that cannot be run on the
available machine within the available time is unanswered in a way that has
nothing to do with the first two.

These three sources are independent. A question may be inexpressible in a
model whose corpus is rich and whose engine is fast; it may be expressible and
supported but unentailed; it may be expressible, entailed and beyond the
budget. Our formal development takes this independence as its starting point,
and \Cref{thm:independence} makes it precise.

\subsection{What follows from taking incompleteness seriously}

If incompleteness is expected, then the informative content of a knowledge
graph is not only the set of questions it answers. It is the \emph{structured
account of the questions it does not answer}. A system that answers
sixty per cent of a question set and cannot characterise the other forty is
less useful than one that answers forty per cent and, for each of the
remaining sixty, states which of the three sources is responsible and what
would have to change to unblock it. The second system supports a decision; the
first supports only a percentage.

This reframing has four consequences that we develop in the body of the paper.

\begin{enumerate}[label=(C\arabic*),leftmargin=*]
\item \textbf{Refusal is a first-class outcome.} A verdict of \vexp{} is a
      result. It is derived, it is reproducible, it constrains future work,
      and it may be cited. Treating refusal as a failure to be minimised makes
      an honest refusal look worse than a confident wrong answer.
      (\Cref{sec:verdicts}.)
\item \textbf{A capability claim describes a compiler, not a formalism.} The
      statement ``this system supports counting'' is true or false of a
      lowering procedure, not of the target logic. The two come apart in both
      directions, and only the compiler-level reading supports a decidable
      containment check. (\Cref{sec:capability}.)
\item \textbf{Every constraint carries a control that can invalidate its own
      result.} A constraint may be present in the compiled form and yet do no
      work. Detecting this requires two checks that do not subsume one
      another. (\Cref{sec:vacuity}.)
\item \textbf{Derivation happens once.} Every view of a result is computed
      from a single derivation, so that two presentations of the same system
      cannot disagree. (\Cref{sec:single}.)
\end{enumerate}

\subsection{Relation to existing work, and what is new}

Open-world semantics and the absence of a unique-name assumption are standard
in description logic~\citep{BaaderCalvanese2017,Horrocks2003}, and the
distinction between an empty extension and an unentailed count is not
new. Benchmarking suites for ontology reasoning report per-engine timeouts and
errors~\citep{Guo2005LUBM,Singh2020OWL2Bench,Parsia2017ORE}, and it is common
practice to report a reasoner's failure to terminate rather than to record a
zero.

What is new here is threefold. First, the separation of \emph{inexpressible in
the logic} from \emph{unaskable through this interface}: benchmark suites
typically report an engine failure, folding together a limitation of a profile
with a limitation of a particular adapter, and the fold systematically
converts an implementation defect into a claim about a standard. Second, the
treatment of \emph{vacuity} as a first-class hazard requiring two irredundant
checks, which to our knowledge has not been formalised for
constraint-lowering pipelines even though the analogous notion of vacuous
satisfaction is well developed in model
checking~\citep{Beer2001,Kupferman2003}. Third, the \emph{non-degeneracy}
requirement---no verdict other than \vans{} may carry an empty answer set---
which turns the conflation described in \S1.1 into a checkable structural
property of a system rather than a matter of reporting discipline.

\subsection{Organisation}

\Cref{sec:setting} fixes the setting and defines models, corpora, engines,
questions and features. \Cref{sec:three} defines expressibility, derivability
and answerability and proves their pairwise independence. \Cref{sec:verdicts}
constructs the verdict space and proves non-degeneracy. \Cref{sec:capability}
develops capability sets as compiler-level claims and proves the containment
theorem and the two failure directions. \Cref{sec:vacuity} treats vacuity and
proves the irredundance of the two checks. \Cref{sec:blockers} gives the
blocker trichotomy and the unblocking relation. \Cref{sec:semantics} treats
the divergence of semantics between target languages, showing that the same
model lowered to two engines yields legitimate disagreement that is a result
rather than a defect. \Cref{sec:single} proves the single-derivation property
and the staleness theorem for status fields. \Cref{sec:validation} states the
seventeen computational propositions and reports their measured outcomes.
\Cref{sec:discussion} discusses limits and what the framework does not do.

\section{Setting}
\label{sec:setting}

We work with four kinds of object. Throughout, all sets are finite unless
stated otherwise.

\begin{definition}[Signature]
\label{def:signature}
A \emph{signature} $\Sigma = (N_C, N_R, N_I)$ consists of disjoint finite sets
of \emph{concept names}, \emph{role names} and \emph{individual names}.
\end{definition}

\begin{definition}[Model]
\label{def:model}
A \emph{model} $\Model = (\Sigma, \mathcal{T})$ is a signature together with a
finite set $\mathcal{T}$ of axioms over $\Sigma$, drawn from a fixed axiom
grammar. We write $\Feat(\Model) \subseteq \Feat$ for the set of
\emph{features} occurring in $\mathcal{T}$, where $\Feat$ is a fixed finite
alphabet of construct names fixed in \Cref{def:features}.
\end{definition}

\begin{definition}[Corpus]
\label{def:corpus}
A \emph{corpus} $\Corpus$ over $\Sigma$ is a finite set of ground assertions
$A(a)$ and $r(a,b)$ with $A \in N_C$, $r \in N_R$, $a,b \in N_I$. The pair
$(\Model, \Corpus)$ is a \emph{knowledge base}.
\end{definition}

We deliberately separate $\Model$ from $\Corpus$. The separation is what makes
the second and third blockers of \Cref{sec:blockers} distinguishable: a
question may be unanswerable because the axioms cannot state it or because the
assertions do not populate it, and merging the two into a single ``ontology''
makes those cases indistinguishable by construction.

\begin{definition}[Question]
\label{def:question}
A \emph{question} $\qq$ is a pair $(\varphi, \bar{x})$ where $\varphi$ is a
formula over $\Sigma$ in a fixed query grammar and $\bar{x}$ is a tuple of its
free variables. We write $\Ques$ for the set of questions. Each question
carries a \emph{feature requirement} $\req(\qq) \subseteq \Feat$: the set of
constructs that any faithful rendering of $\qq$ must use.
\end{definition}

\begin{definition}[Feature alphabet]
\label{def:features}
$\Feat$ is the finite alphabet
\[
\Feat = \{\,
\textsf{SUBSUMPTION},\;
\textsf{EXISTENTIAL},\;
\textsf{VALUE},\;
\textsf{NEGATION},\;
\textsf{COUNTING},\;
\textsf{TRANSITIVE\_AXIOM},\;
\textsf{TRANSITIVE\_QUERY},\;
\textsf{BOUNDED\_PATH},
\]
\[
\textsf{DISJOINTNESS},\;
\textsf{UNIQUE\_NAMES}
\,\}.
\]
\end{definition}

Two remarks on this alphabet are load-bearing and are proved later.

\begin{remark}[Splitting one name into two]
\label{rem:transitive-split}
Transitivity appears twice. \textsf{TRANSITIVE\_AXIOM} is the assertion that a
role $r$ is transitive, discharged by the reasoner when it closes the
interpretation. \textsf{TRANSITIVE\_QUERY} is the requirement that a query
compute the transitive closure of a role at answering time. These are distinct
mechanisms: a system may implement either without the other. Naming them
identically makes a containment check unsound in a way \Cref{prop:half-true}
makes precise.
\end{remark}

\begin{remark}[Listing what you do lower]
\label{rem:listing}
\textsf{DISJOINTNESS} and \textsf{UNIQUE\_NAMES} are features of a
\emph{model}, not of a question, and it is tempting to omit them from an
alphabet indexed by questions. Omitting them is what makes a defect
undetectable: if the alphabet cannot name a construct, then a compiler that
silently ignores that construct in $\mathcal{T}$ produces a knowledge base
differing from its sibling's, and any comparison between the two engines
attributes the difference to the engines. \Cref{prop:nameable} states this.
\end{remark}

\begin{definition}[Engine]
\label{def:engine}
An \emph{engine} $\Engine = (\lowr_\Engine, \sup(\Engine), \mathrm{solve}_\Engine)$
consists of a partial \emph{lowering} map $\lowr_\Engine$ taking a model and a
question to a program in the engine's target language, a \emph{declared
capability set} $\sup(\Engine) \subseteq \Feat$, and a \emph{solver}
$\mathrm{solve}_\Engine$ taking a program and a corpus to an answer set
together with a certification flag and an elapsed time.
\end{definition}

\begin{assumption}[Determinism and budget]
\label{ass:budget}
$\mathrm{solve}_\Engine$ is deterministic given $(\text{program}, \Corpus)$
and is run under an explicit wall-clock budget $\tau$. Any run exceeding $\tau$
is terminated and produces no answer set.
\end{assumption}

\Cref{ass:budget} is not a technicality. It is what makes \vtime{} a verdict
rather than a crash: a terminated run that returned an empty answer set would
be recorded by any differential comparison as \emph{agreement} between a
working engine and a hung one.

\section{Three independent predicates}
\label{sec:three}

\subsection{Definitions}

\begin{definition}[Expressible]
\label{def:expressible}
A question $\qq$ is \emph{expressible in} $\Model$, written
$\mathrm{Exp}(\qq, \Model)$, if the constructs $\req(\qq)$ are all available in
the axiom and query grammar over $\Sigma(\Model)$, and every name occurring in
$\qq$ is in $\Sigma(\Model)$.
\end{definition}

Expressibility is a property of the \emph{language}. It does not depend on the
corpus, and it does not depend on any engine.

\begin{definition}[Derivable]
\label{def:derivable}
Let $\qq = (\varphi, \bar{x})$ be expressible in $\Model$. A tuple $\bar{a}$ of
individuals is a \emph{certain answer} if
$(\Model, \Corpus) \models \varphi[\bar{a}/\bar{x}]$, i.e.\ $\varphi[\bar{a}]$
holds in every model of the knowledge base. The question is \emph{derivable},
written $\mathrm{Der}(\qq, \Model, \Corpus)$, if the answer to $\qq$ is
determined by the knowledge base: for every tuple $\bar{a}$ over $N_I$, either
$(\Model, \Corpus) \models \varphi[\bar{a}]$ or
$(\Model, \Corpus) \models \neg\varphi[\bar{a}]$.
\end{definition}

Derivability is a property of the \emph{semantics}. Note carefully that it is
stronger than the existence of certain answers: a question with no certain
answers may still be underdetermined, and the difference between ``the
extension is empty'' and ``the extension is not determined'' is exactly the
one \S1.1 identified as invisible.

\begin{definition}[Answerable]
\label{def:answerable}
A question $\qq$ is \emph{answerable by} $\Engine$ \emph{on}
$(\Model, \Corpus)$ \emph{within} $\tau$, written
$\mathrm{Ans}_\tau(\qq, \Model, \Corpus, \Engine)$, if $\lowr_\Engine(\Model, \qq)$
is defined and $\mathrm{solve}_\Engine$ on the resulting program and $\Corpus$
terminates within $\tau$ with certification.
\end{definition}

Answerability is a property of the \emph{implementation and the budget}. It
mentions a clock. Neither of the other two does.

\subsection{Independence}

\begin{theorem}[Pairwise independence]
\label{thm:independence}
None of $\mathrm{Exp}$, $\mathrm{Der}$, $\mathrm{Ans}_\tau$ is entailed by the
conjunction of the other two. Formally, for each of the three predicates there
exist a model, a corpus, an engine, a budget and a question witnessing the
other two and failing that one.
\end{theorem}

\begin{proof}
We exhibit three witnesses.

\emph{$\mathrm{Der}$ and $\mathrm{Ans}_\tau$ without $\mathrm{Exp}$.} Let
$\Model_1$ have grammar without cardinality restrictions, so
$\textsf{COUNTING} \notin \Feat(\Model_1)$ and no question requiring it is
expressible. Take $\qq_1 = $ ``entities with exactly four $r$-successors''.
Now consider the surrogate question $\qq_1' = $ ``entities $x$ with
$r(x,a_1),\dots,r(x,a_4)$ for the four named individuals $a_1,\dots,a_4$'',
which \emph{is} expressible and which, on a corpus asserting exactly those four
edges for one individual and no others, is derivable and answerable by any
engine handling conjunctive queries. The pair witnesses that expressibility
fails on $\qq_1$ while the semantic and computational conditions are
unproblematic for the surrogate; formally, $\mathrm{Exp}(\qq_1, \Model_1)$ is
false while for the extension $\Model_1^+$ obtained by adding cardinality to
the grammar, $\mathrm{Der}(\qq_1, \Model_1^+, \Corpus)$ and
$\mathrm{Ans}_\tau(\qq_1, \Model_1^+, \Corpus, \Engine)$ both hold. Since
$\Corpus$ and $\Engine$ are unchanged, the failure is attributable to the
grammar alone.

\emph{$\mathrm{Exp}$ and $\mathrm{Ans}_\tau$ without $\mathrm{Der}$.} Let
$\Model_2$ have cardinality restrictions, and let $\Corpus_2$ assert
$r(c,a_1),\dots,r(c,a_4)$ with no disjointness among the $a_i$ and no closure
axiom on $r$. Then $\qq_1$ is expressible, and an engine returns promptly. But
$(\Model_2, \Corpus_2) \not\models \text{``$c$ has exactly four
$r$-successors''}$, because there is a model of the knowledge base in which
$a_1 = a_2$ (no unique-name assumption) and another in which $c$ has a fifth
successor (no closure). Nor does the knowledge base entail the negation, since
a model realising exactly four distinct successors exists. So the question is
underdetermined: $\mathrm{Der}$ fails while the other two hold.

\emph{$\mathrm{Exp}$ and $\mathrm{Der}$ without $\mathrm{Ans}_\tau$.} Let
$\Model_3$ and $\Corpus_3$ be such that $\qq$ is expressible and determined---
for instance, take $\Corpus_3$ to assert pairwise distinctness of the $a_i$
and a closure axiom on $r$ at $c$, which fixes the count---and let $\tau$ be
smaller than the time $\mathrm{solve}_\Engine$ requires on
$\lowr_\Engine(\Model_3,\qq)$. Such a $\tau$ exists because the solver's
running time on a fixed input is a positive real. Then $\mathrm{Ans}_\tau$
fails while the other two hold. Alternatively, and without appeal to the
budget, let $\lowr_\Engine$ be undefined on $\qq$---the engine's adapter has no
surface for cardinality---which also falsifies $\mathrm{Ans}_\tau$ while
leaving $\mathrm{Exp}$ and $\mathrm{Der}$ untouched.
\end{proof}

\begin{corollary}[The conflation]
\label{cor:conflation}
A system whose output on a question is a single answer set cannot distinguish
$\neg\mathrm{Exp}$, $\neg\mathrm{Der}$ and $\neg\mathrm{Ans}_\tau$, since all
three admit the empty set as output and, by \Cref{thm:independence}, no two of
them determine the third.
\end{corollary}

\begin{remark}[The two witnesses of $\neg\mathrm{Ans}_\tau$ are different]
\label{rem:two-ans}
The proof gave two independent ways to falsify answerability: the budget and
the undefined lowering. They warrant different verdicts and different actions,
and \Cref{sec:verdicts} separates them as \vtime{} and \vsurf. Merging them is
how ``this adapter has no cardinality surface'' becomes ``this profile cannot
count''.
\end{remark}

\subsection{The specific structure of counting}

The second witness in \Cref{thm:independence} deserves a sharper statement,
because it is the case most often mishandled.

\begin{proposition}[Two requirements, one aggregate]
\label{prop:counting-split}
Let $\qq_{\geq n}$, $\qq_{\leq n}$ and $\qq_{=n}$ be the questions
``$x$ has at least / at most / exactly $n$ $r$-successors''. Under open-world
semantics without a unique-name assumption:
\begin{enumerate}[label=(\roman*)]
\item $\qq_{\geq n}$ is derivable only if the counted successors are known
      pairwise distinct;
\item $\qq_{\leq n}$ is derivable only if the successor set is closed;
\item $\qq_{=n}$ requires both.
\end{enumerate}
\end{proposition}

\begin{proof}
(i) Suppose $r(c, a_1), \dots, r(c, a_n)$ hold and distinctness is not
asserted. The interpretation identifying all $a_i$ satisfies the knowledge
base and gives $c$ one successor, so ``at least $n$'' is not entailed for
$n>1$. Conversely if the $a_i$ are pairwise distinct in every model then every
model gives $c$ at least $n$ successors. (ii) Suppose the successor set is not
closed. Extend any model by a fresh element $d$ with $r(c,d)$; if the extension
satisfies the knowledge base, ``at most $n$'' fails in it. Conversely a closure
axiom $c \sqsubseteq \forall r.\{a_1,\dots,a_n\}$ makes every model's successor
set a subset of $\{a_1^\mathcal{I},\dots,a_n^\mathcal{I}\}$, bounding the count
by $n$. (iii) $\qq_{=n}$ is the conjunction, and a conjunction is entailed only
if both conjuncts are.
\end{proof}

\begin{remark}[The tempting check is wrong in the hiding direction]
\label{rem:una-scope}
A system implementing \Cref{prop:counting-split}(i) must check that the
\emph{counted individuals} are known distinct. It is tempting instead to check
whether the model contains \emph{any} distinctness assertion, because that is
a single boolean over $\mathcal{T}$. This is unsound in the direction that
conceals the problem: a model may assert distinctness over an unrelated pair of
individuals---two role constants, say---while saying nothing about the
successors being counted, and the boolean is then true while the requirement of
(i) is unmet. A distinctness assertion somewhere in a model is not a
distinctness assertion over the counted set.
\end{remark}

\begin{remark}[False by construction, not by inspection]
\label{rem:by-construction}
If a model's grammar contains no closure construct then
\Cref{prop:counting-split}(ii) fails for structural reasons, and a system
should record this fact as \emph{holding by construction}. The distinction
matters for maintenance: a fact recorded as ``false by construction'' is
flipped by extending the grammar, whereas a fact recorded as ``false by
inspection'' invites someone to look for a closure axiom that the grammar
cannot express. In the first framing, extending the grammar is what changes the
verdict; in the second, a test passing is.
\end{remark}

\section{The verdict space}
\label{sec:verdicts}

\subsection{Construction}

\begin{definition}[Verdict]
\label{def:verdict}
A \emph{verdict} is a pair $\Verd = (\ell, \pi)$ where $\ell$ is a label from
\[
L = \{\vans,\ \vexp,\ \vder,\ \vsurf,\ \vtime,\ \vvac,\ \vinert\}
\]
and $\pi$ is a label-specific \emph{payload}. The payloads are:
\begin{center}
\begin{tabular}{@{}ll@{}}
\toprule
Label & Payload \\
\midrule
\vans   & answer set $\Ans$, certification flag, elapsed time \\
\vexp   & the subset $\req(\qq) \setminus \Feat(\Model)$ responsible \\
\vder   & which clause of \Cref{prop:counting-split} (or analogue) failed \\
\vsurf  & the feature the lowering has no surface for \\
\vtime  & the budget $\tau$ that was exceeded \\
\vvac   & the compiled expressions found identical \\
\vinert & the two answer sets found identical \\
\bottomrule
\end{tabular}
\end{center}
\end{definition}

\begin{definition}[Evaluation]
\label{def:evaluation}
Given $(\Model, \Corpus, \Engine, \tau)$ and a question $\qq$, the
\emph{evaluation} $\mathrm{Ev}(\qq)$ is computed in the following order, the
first applicable clause determining the verdict:
\begin{enumerate}[label=(E\arabic*),leftmargin=*]
\item if $\req(\qq) \not\subseteq \Feat(\Model)$, return
      $(\vexp, \req(\qq)\setminus\Feat(\Model))$;
\item if $\qq$ is not derivable by the applicable clause of
      \Cref{prop:counting-split}, return $(\vder, \text{clause})$;
\item if $\req(\qq) \not\subseteq \sup(\Engine)$, return
      $(\vsurf, \req(\qq)\setminus\sup(\Engine))$;
\item compile $P = \lowr_\Engine(\Model,\qq)$ and, if $\qq$ carries a
      droppable constraint, the control $P^\circ$ obtained by removing it. If
      $P = P^\circ$ syntactically, return $(\vvac, (P,P^\circ))$;
\item run $P$ within $\tau$; if \emph{it} exceeds $\tau$, return
      $(\vtime, \tau)$. Run $P^\circ$ within $\tau$; if the control alone
      exceeds it, the vacuity test of (E6) is \emph{inconclusive} and is
      recorded as such in the payload of whatever clause below applies;
\item if $\Ans(P) = \Ans(P^\circ)$ \emph{and} $\Ans(P) \neq \emptyset$, return
      $(\vinert, \Ans(P))$;
\item otherwise return $(\vans, (\Ans(P), \text{certified}, t))$.
\end{enumerate}
\end{definition}

The order is not arbitrary. Each clause is checked before any clause whose
evidence it would corrupt. In particular (E1) precedes (E3), so that a
limitation of the language is never reported as a limitation of the adapter;
and (E4) precedes (E5), so that a vacuous control is detected before a budget
is spent proving that it agrees.

Two side conditions in (E5) and (E6) deserve comment, because the natural
statement of each is wrong in a way that only shows up when the clause fires.

\begin{remark}[The control's budget is not the question's]
\label{rem:control-budget}
A control that exceeds $\tau$ is not a \vtime{} verdict for $\qq$. The question
has been answered; what has failed is the auxiliary test of whether its
constraint did any work. Returning \vtime{} here would report a property of
$P^\circ$ as a verdict on $\qq$ --- precisely the conflation of
\Cref{cor:conflation}, one level down. The case is not exotic: whenever the
droppable constraint is the one that made the question affordable, the control
is by construction the more expensive program, so this is the \emph{expected}
behaviour for exactly the questions of \Cref{prop:bounded-cheaper} rather than
a rare degeneracy. An \vans{} verdict whose vacuity test did not complete must
carry that fact in its payload, since otherwise \vans{} silently implies a
check that was never performed.
\end{remark}

\begin{remark}[Two empty sets agree about nothing]
\label{rem:empty-agreement}
The non-emptiness side condition in (E6) is what prevents a capability defect
from being reported as a vacuity finding. \Cref{def:sem-vacuity} infers that a
constraint did no work from the observation that two \emph{differing} programs
produce the same answers, and that inference requires the agreement to have
content. If neither program grounds --- the ungrounded-aggregate case of
\Cref{prop:under} --- both return $\emptyset$ and agree vacuously, for a reason
having nothing to do with the constraint. Without the side condition, a
compiler emitting a program that never runs is labelled \vinert, which is a
verdict about the question rather than about the compiler, and the defect is
concealed by the very machinery built to expose it.
\end{remark}

\begin{proposition}[Order soundness]
\label{prop:order}
If $\mathrm{Ev}(\qq) = (\ell, \pi)$ with $\ell \neq \vans$ then every clause
preceding $\ell$'s clause in \Cref{def:evaluation} is inapplicable, and the
payload $\pi$ names a witness sufficient to reproduce the verdict without
re-running the engine.
\end{proposition}

\begin{proof}
Immediate from the first-applicable rule. For the second part, inspect the
payload column of \Cref{def:verdict}: each payload is a finite object computed
from $\Model$, $\qq$, $\sup(\Engine)$, $\lowr_\Engine$ and $\tau$ alone, none
of which depends on the solver's execution except in the \vans{} and \vinert{}
rows, whose payloads include the answer sets themselves.
\end{proof}

\subsection{Non-degeneracy}

The property that makes the verdict space more than a taxonomy is the
following.

\begin{theorem}[Non-degeneracy]
\label{thm:nondegeneracy}
Under \Cref{def:evaluation}, no verdict with $\ell \neq \vans$ carries an
empty answer set, because no such verdict carries an answer set at all.
Consequently the observation ``the output was empty'' determines
$\ell = \vans$, and the five situations of \S1.1 are pairwise distinguishable
from the verdict alone.
\end{theorem}

\begin{proof}
By inspection of the payload table: the only labels whose payload contains an
answer set are \vans{} and \vinert. \vinert{} is reached only when
$\Ans(P) = \Ans(P^\circ)$, and its payload records that common set as evidence
of inertness rather than as an answer, so no consumer reading the label as an
answer set is licensed. For the second claim, suppose a consumer observes an
empty output. By the first part the label must be \vans. The five situations
of \S1.1 are then separated as follows: an inexpressible question yields
\vexp{} by (E1); an unentailed count yields \vder{} by (E2); a genuinely empty
extension yields $(\vans, (\emptyset, \text{certified}, t))$; a dropped
constraint yields \vvac{} by (E4) if the drop is syntactic and \vinert{} by
(E6) if it is semantic; an exhausted budget yields \vtime{} by (E5). These are
five distinct labels.
\end{proof}

\begin{corollary}[Certified emptiness is a result]
\label{cor:certified-empty}
A verdict $(\vans, (\emptyset, \text{certified}, t))$ is a positive claim:
the extension is empty and the run that established this was complete. It is
therefore not interchangeable with any $\ell \neq \vans$, and in particular
not with \vtime, whose payload asserts nothing about the extension.
\end{corollary}

\begin{remark}[Certification is reported, never folded in]
\label{rem:certification}
The certification flag is a component of the \vans{} payload and not a filter
on it. A system that dropped uncertified answers would report a search-limited
run and a complete run identically when both return the same set, and would
report a search-limited run and an exhausted one identically when the former
returns nothing. Where a solver mode is constitutionally unable to certify,
the flag must be explicitly false rather than absent: an unset field reads as
a default and a default reads as satisfactory.
\end{remark}

\begin{remark}[The budget is part of the claim]
\label{rem:budget-in-payload}
$(\vtime, 120\,\mathrm{s})$ and $(\vtime, 3600\,\mathrm{s})$ are different
claims, and a verdict space that records only the label loses the difference.
This is why \Cref{def:verdict} makes $\tau$ the payload rather than metadata.
\end{remark}

\subsection{Refusal as a measurement}

\begin{proposition}[Comparability of refusals]
\label{prop:compare-refusals}
Let $\Engine_1, \Engine_2$ be engines evaluated on the same
$(\Model,\Corpus,\qq,\tau)$. Then the pair
$(\mathrm{Ev}_1(\qq), \mathrm{Ev}_2(\qq))$ is informative even when neither
label is \vans, and in particular the pair $(\vexp, \vans)$ is not evidence
that $\Engine_2$ is more capable than $\Engine_1$ on $\qq$ unless
$\Engine_2$'s payload includes certification.
\end{proposition}

\begin{proof}
Suppose $\mathrm{Ev}_1(\qq) = (\vexp, \{\textsf{COUNTING}\})$ and
$\mathrm{Ev}_2(\qq) = (\vans, (\emptyset, \cdot, t))$ on a corpus whose ground
truth for $\qq$ is a non-empty set of size $m > 0$. Then $\Engine_1$'s output
correctly identifies that $\qq$ cannot be posed in its language, while
$\Engine_2$'s output differs from ground truth by $m$. Ranking $\Engine_2$
above $\Engine_1$ on the basis that it produced an answer set orders them by
output type rather than by correctness, and reverses the correct ordering
whenever $m>0$. If $\Engine_2$'s empty set is certified, then by
\Cref{cor:certified-empty} its verdict is a positive claim of emptiness, which
contradicts $m>0$ and identifies a defect; if it is uncertified, the verdict
carries no claim about the extension at all.
\end{proof}

\Cref{prop:compare-refusals} is the formal statement of the phenomenon that
motivated this paper. An engine that refuses honestly appears \emph{less}
capable than one returning a confident wrong zero, under any scoring rule that
counts answers rather than verdicts.

\section{Capability is a property of the compiler}
\label{sec:capability}

\subsection{The containment check}

\begin{definition}[Faithful lowering]
\label{def:faithful}
$\lowr_\Engine$ is \emph{faithful for} $f \in \Feat$ if for every model and
question in which $f$ occurs, the program $\lowr_\Engine(\Model,\qq)$ has the
same certain answers as $\qq$ has under \Cref{def:derivable}.
\end{definition}

\begin{definition}[Honest capability set]
\label{def:honest}
$\sup(\Engine)$ is \emph{honest} if $f \in \sup(\Engine)$ iff
$\lowr_\Engine$ is faithful for $f$.
\end{definition}

\begin{theorem}[Containment]
\label{thm:containment}
Let $\sup(\Engine)$ be honest. If $\req(\qq) \cup \Feat(\Model) \subseteq
\sup(\Engine)$ then $\lowr_\Engine(\Model,\qq)$ is defined and faithful. If
$\req(\qq) \cup \Feat(\Model) \not\subseteq \sup(\Engine)$ then
\Cref{def:evaluation}(E3) fires and the verdict is $(\vsurf, \cdot)$, which by
\Cref{thm:nondegeneracy} carries no answer set.
\end{theorem}

\begin{proof}
The first part is \Cref{def:honest} applied featurewise, using that
faithfulness is closed under the conjunction of independent constructs by
\Cref{def:faithful}. The second is by construction of (E3), and the absence of
an answer set is \Cref{thm:nondegeneracy}.
\end{proof}

The theorem is only as strong as its hypothesis, and the hypothesis is what
the next two subsections attack.

\subsection{The formalism can do what the compiler cannot}

\begin{proposition}[Under-declaration is correct, not modest]
\label{prop:under}
There exist a feature $f$, an engine $\Engine$ and a target language $T$ such
that $T$ supports $f$ in full generality, $\lowr_\Engine$ produces a program
in $T$ using $f$, that program is rejected or fails to certify, and honesty of
$\sup(\Engine)$ therefore requires $f \notin \sup(\Engine)$.
\end{proposition}

\begin{proof}
Take $f = \textsf{COUNTING}$ and $T$ a rule language with aggregate support.
Counting is expressible in $T$ in full generality. Let $\lowr_\Engine$ emit,
for the question ``entities with exactly four $r$-successors'', a rule whose
body pairs an aggregate over an unbound successor variable with an equality
test on the aggregate's result. If the engine certifies answers only within a
fragment for which grounding is bounded, this program falls outside that
fragment: the aggregate's grounding is unbounded in the successor variable and
the engine rejects it. Removing the fragment gate does not repair the
situation---it produces a run that terminates with an empty answer set, which
is the answer set of a program whose aggregate never grounded, and which by
\Cref{cor:certified-empty} would be read as a positive claim of emptiness.
Since $\lowr_\Engine$ is faithful for $f$ in neither configuration,
\Cref{def:honest} forces $f \notin \sup(\Engine)$, notwithstanding that $T$
counts perfectly well.
\end{proof}

\begin{corollary}[A capability recommendation must be tested by attempting it]
\label{cor:test-the-recommendation}
An audit of $\sup(\Engine)$ conducted against the target language $T$ rather
than against $\lowr_\Engine$ will recommend adding features that must not be
added. Such a recommendation is refuted only by writing the lowering and
observing its verdict.
\end{corollary}

\begin{remark}[One set cannot carry two positions]
\label{rem:one-set}
$\sup(\Engine)$ as defined is a single set checked against both
$\Feat(\Model)$ and $\req(\qq)$. It therefore cannot express ``supported when
the construct occurs in a query, unsupported when it occurs in an axiom'',
which is a real distinction: lowering a construct into a rule body and
lowering it into a rule head are different problems with different fragment
conditions. A system needing that distinction requires two sets, and a system
with one set must record which position its declaration refers to.
\end{remark}

\subsection{One name for two mechanisms}

\begin{proposition}[Half-true containment]
\label{prop:half-true}
Suppose $\Feat$ contains a single symbol $\textsf{TRANSITIVE}$ conflating
\Cref{rem:transitive-split}'s two mechanisms. Then there exists an engine
$\Engine$ for which $\textsf{TRANSITIVE} \in \sup(\Engine)$ is neither true
nor false, in the sense that \Cref{def:honest} is unsatisfiable: including it
makes $\sup$ dishonest for questions requiring the query mechanism, and
excluding it makes $\sup$ dishonest for models using the axiom mechanism.
\end{proposition}

\begin{proof}
Let $\Engine$ lower to a description logic supporting transitive role axioms
whose query language has no closure operator. Then $\lowr_\Engine$ is
faithful for \textsf{TRANSITIVE\_AXIOM} and unfaithful for
\textsf{TRANSITIVE\_QUERY}. Under \Cref{def:honest} with the split alphabet,
$\sup(\Engine)$ contains the first and not the second, and
\Cref{thm:containment} applies. Under the conflated alphabet, honesty demands
$\textsf{TRANSITIVE} \in \sup(\Engine)$ (from faithfulness on axioms) and
$\textsf{TRANSITIVE} \notin \sup(\Engine)$ (from unfaithfulness on queries)
simultaneously.
\end{proof}

\begin{remark}[Do not make the base relation transitive]
\label{rem:no-base-transitive}
A tempting repair for a missing query-side closure is to declare the base role
transitive in $\mathcal{T}$. This is not a repair but a change of model:
declaring a participation role transitive entails that an entity participates
in the participants of its participants, corrupting every other question over
the same knowledge base. The correct construction introduces a fresh
super-role, asserts \emph{it} transitive, and answers the closure question over
the super-role, leaving the base role's extension untouched.
\end{remark}

\subsection{Nameability}

\begin{proposition}[An unnameable feature is an undetectable omission]
\label{prop:nameable}
Let $f$ be a feature of models that is absent from $\Feat$. Let $\Engine_1$
lower $f$ faithfully and $\Engine_2$ ignore it. Then for any question,
\Cref{def:evaluation} never reaches a verdict attributing a discrepancy to
$f$, and any differential comparison of $\Engine_1$ and $\Engine_2$ compares
two different knowledge bases while reporting a difference between two
engines.
\end{proposition}

\begin{proof}
(E1) and (E3) are containment tests over $\Feat$; a symbol absent from $\Feat$
occurs in neither $\req(\qq)$ nor $\Feat(\Model)$ nor $\sup(\Engine)$, so
neither test can fire on it. The evaluations therefore proceed to (E4)--(E7)
in both engines. $\Engine_2$'s program is compiled from $\mathcal{T}$ minus the
axioms mentioning $f$, so its answer sets are those of a knowledge base
differing from $\Engine_1$'s. The reported difference is labelled by engine
because the engine is the only varying parameter the evaluation records.
\end{proof}

\Cref{prop:nameable} is the formal content of \Cref{rem:listing}: enumerating
the constructs one \emph{does} lower is what makes an omission on the other
side detectable, and the alphabet must therefore be extended whenever a new
construct enters any model, not only when a question requires it.

\section{Vacuity}
\label{sec:vacuity}

A constraint may survive lowering in form and not in effect. This section
formalises the two ways that happens and proves that catching them requires
two checks.

\begin{definition}[Control]
\label{def:control}
Let $\qq$ carry a \emph{droppable constraint} $c$---a negation, a cardinality
restriction, or a path bound. The \emph{control} $\qq^\circ$ is $\qq$ with $c$
removed. The compiled pair is $P = \lowr_\Engine(\Model,\qq)$,
$P^\circ = \lowr_\Engine(\Model,\qq^\circ)$.
\end{definition}

\begin{definition}[Syntactic vacuity]
\label{def:syn-vacuity}
$c$ is \emph{syntactically vacuous} at $\Engine$ if $P = P^\circ$ as
programs.
\end{definition}

\begin{definition}[Semantic vacuity]
\label{def:sem-vacuity}
$c$ is \emph{semantically vacuous} at $(\Engine,\Corpus)$ if $P \neq P^\circ$
but $\Ans(P) = \Ans(P^\circ)$ on $\Corpus$.
\end{definition}

\begin{theorem}[Irredundance]
\label{thm:irredundance}
Neither check subsumes the other:
\begin{enumerate}[label=(\roman*)]
\item there exist $(\Engine, \Corpus, \qq, c)$ with $c$ syntactically vacuous
      and $\Ans(P) = \Ans(P^\circ)$ trivially, which the semantic check
      reports but with a payload identifying the wrong cause;
\item there exist $(\Engine, \Corpus, \qq, c)$ with $c$ syntactically present,
      $P \neq P^\circ$, and $\Ans(P) = \Ans(P^\circ)$, which the syntactic
      check cannot detect.
\end{enumerate}
\end{theorem}

\begin{proof}
(ii) is the substantive direction. Let $c$ be a negated conjunction
constraining a \emph{role filler}---``$x$ has no $r$-successor that is a
$B$''---and let $\lowr_\Engine$ construct the negated expression by rooting
each negated block at the block's own first variable. The lowered expression
then reads ``$x$ is not itself a $B$-with-role'' rather than ``$x$ has no
$B$-successor'': it is a different, well-formed expression, so $P \neq P^\circ$
and the syntactic check passes. If in addition the model asserts pairwise
disjointness between the class of $x$ and the class $B$, the mis-rooted
expression is entailed of every $x$ and the constraint is satisfied
identically, giving $\Ans(P) = \Ans(P^\circ)$. Thus a semantic check is
required, and its payload---the two identical answer sets---names the actual
symptom.

(i) Let $\lowr_\Engine$ ignore $c$ entirely, so that $P$ and $P^\circ$ are
literally the same program. Then $\Ans(P) = \Ans(P^\circ)$ holds trivially and
the semantic check fires, but its payload asserts that the constraint made no
difference to the answers, when the stronger and more actionable fact is that
the constraint never reached the program. A syntactic check reports that fact
directly, and reports it without spending the budget of two solver runs.
\end{proof}

\begin{corollary}[A control must differ from what it controls]
\label{cor:control-differs}
$P \neq P^\circ$ is a precondition for $P^\circ$ to function as a control. A
pipeline that constructs the control by textual manipulation of $\qq$ and
recompiles may produce $P^\circ = P$ without any error being raised, in which
case the subsequent comparison is between a program and itself and can only
report agreement.
\end{corollary}

\begin{remark}[Trigger on droppability, not on a construct name]
\label{rem:trigger}
The set of questions carrying a control must be defined by the property
``carries a constraint whose lowering could be silently dropped'' rather than
by naming one construct such as negation. Cardinality restrictions and path
bounds are droppable by the same mechanism---each is a modifier on an
otherwise well-formed pattern, so removing it leaves a program that still
compiles---and a trigger enumerating only negation admits them silently.
\end{remark}

\begin{proposition}[Vacuity defeats agreement as evidence]
\label{prop:agreement}
Let $\Engine_1, \Engine_2$ agree on $\qq$, i.e.\
$\Ans_1(P_1) = \Ans_2(P_2)$. If $c$ is vacuous at $\Engine_1$ in either sense,
the agreement is not evidence that either lowering is faithful for $c$.
\end{proposition}

\begin{proof}
By \Cref{def:sem-vacuity}, $\Ans_1(P_1) = \Ans_1(P_1^\circ)$, so
$\Engine_1$'s output is the output of the unconstrained question. Agreement
with $\Engine_2$ therefore establishes at most that $\Engine_2$ agrees with
$\Engine_1$'s \emph{unconstrained} answer, which is consistent with
$\Engine_2$ also dropping $c$ and equally consistent with $c$ being
extensionally inert on this corpus. No conclusion about faithfulness for $c$
follows.
\end{proof}

\section{The blocker trichotomy}
\label{sec:blockers}

We now attach to each unanswered question a statement of what would have to
change.

\begin{definition}[Blocker]
\label{def:blocker}
For a question $\qq$ unanswered under $(\Model,\Corpus,\Engine,\tau)$, its
\emph{blocker} is:
\begin{itemize}[leftmargin=*]
\item \textsf{blocked-by-model} if $\neg\mathrm{Exp}(\qq,\Model)$, or if
      $\mathrm{Exp}$ holds but $\neg\mathrm{Der}$ for a reason attributable to
      a missing axiom (absent closure, absent distinctness);
\item \textsf{blocked-by-corpus} if $\mathrm{Exp}$ and the axioms suffice but
      $\Corpus$ contains no assertions populating the question's pattern;
\item \textsf{blocked-by-engine} if $\mathrm{Exp}$ and $\mathrm{Der}$ hold but
      $\neg\mathrm{Ans}_\tau$.
\end{itemize}
\end{definition}

\begin{definition}[Unblocking]
\label{def:unblocking}
An \emph{unblocker} for a blocked question is a change $\delta$ to exactly one
of $\Model$, $\Corpus$, $\Engine$ or $\tau$ such that the question is
answerable after $\delta$. A blocker record is \emph{actionable} if it names
at least one unblocker.
\end{definition}

\begin{proposition}[Blockers are engine-indexed or universally quantified]
\label{prop:engine-index}
A \textsf{blocked-by-model} or \textsf{blocked-by-corpus} record carrying no
engine is a claim about all engines. Consequently a measured \vans{} for the
same question and model, by any engine, is a counterexample to that record and
not an independent observation compatible with it.
\end{proposition}

\begin{proof}
By \Cref{def:blocker}, the first two blockers are conditions on
$(\qq,\Model,\Corpus)$ alone; asserting one without an engine index therefore
asserts it for every $\Engine$. A measured $(\vans,\cdot)$ instantiates the
negation at one $\Engine$, falsifying the universal.
\end{proof}

\begin{corollary}[Superseding, not accumulating]
\label{cor:supersede}
A record set containing both a universally quantified blocker and a measured
\vans{} for the same question and model asserts a claim and its negation. A
system that appends measurements alongside predictions rather than superseding
them will hold such a set while every structural well-formedness check passes,
because well-formedness checks constrain shape and this defect is in content.
\end{corollary}

\begin{proposition}[A measured claim must agree with its evidence about the model]
\label{prop:model-agreement}
Let a record assert $(\vans,\cdot)$ for $\qq$ under model $\Model$, citing an
experiment $e$. If $e$ was executed under $\Model' \neq \Model$, the record is
unsupported regardless of the truth of its content.
\end{proposition}

\begin{proof}
$\mathrm{Ans}_\tau$ is a four-place predicate including the model. An
observation at $\Model'$ instantiates the predicate at $\Model'$ and is silent
at $\Model$.
\end{proof}

\begin{remark}[Which of the two entries is wrong]
\label{rem:which-wrong}
When \Cref{prop:model-agreement} fires against a record that contradicts a
standing prediction, the natural reading is that the prediction is stale. It
may instead be that the prediction is correct and the measured record
misattributes its model. Repairing the contradiction by retiring the
prediction then deletes the true claim and retains the false one. The check
must therefore report which entry failed the agreement test, not merely that a
pair disagreed.
\end{remark}

\begin{remark}[Why only the model is compared]
\label{rem:only-model}
A comparison of the engine field is the obvious companion check and is not
available when an experiment record carries a scalar engine while differential
experiments run two. The asymmetry should be recorded rather than papered
over: a check that silently compares the first engine of a pair reports
agreement on evidence it did not examine.
\end{remark}

\section{Divergent semantics as a result}
\label{sec:semantics}

Lowering one model to two target languages does not in general preserve
meaning, and the residue is informative.

\begin{proposition}[Disjointness lowers to a detector]
\label{prop:disjoint-detector}
Let $\mathcal{T}$ assert that concepts $A,B$ are disjoint, and let a rule
engine lower this as
\[
\mathit{viol}(x) \leftarrow A(x),\, B(x).
\]
This is faithful as a \emph{detector} and unfaithful as a \emph{semantics}: in
a description logic an individual in $A \sqcap B$ renders the knowledge base
inconsistent, from which everything follows, while in the rule engine it
derives exactly one additional fact.
\end{proposition}

\begin{proof}
The rule is pure positive Datalog and its least model contains
$\mathit{viol}(a)$ exactly for those $a$ in both extensions; no other atom is
affected. Under the description-logic reading, $A \sqcap B \sqsubseteq \bot$
together with $A(a), B(a)$ has no model, so every formula is entailed.
\end{proof}

\begin{corollary}[Localisation]
\label{cor:localisation}
On a corpus containing a genuine violation, the description-logic engine
reports global inconsistency, which identifies that something is wrong without
identifying where; the rule engine returns a fact naming the two concepts and
the individual. The verdicts differ, both are correct under their respective
semantics, and only one is diagnostic.
\end{corollary}

\begin{proposition}[Inherited unique names]
\label{prop:inherited-una}
A rule engine over a Herbrand universe satisfies the unique-name assumption
unconditionally, since distinct ground terms denote distinct elements. Hence
\textsf{UNIQUE\_NAMES} is not lowered but inherited, and by
\Cref{prop:counting-split}(i) the two engines legitimately disagree on
$\qq_{\geq n}$: the rule engine derives the lower bound the description logic
declines to entail.
\end{proposition}

\begin{proof}
Herbrand interpretations interpret each constant as itself; distinct constants
are therefore distinct elements, which is the unique-name assumption. Applying
\Cref{prop:counting-split}(i), the distinctness requirement is met in the rule
engine and unmet in the description logic absent explicit assertions.
\end{proof}

\begin{corollary}[The disagreement is the result]
\label{cor:disagree-result}
Under \Cref{prop:inherited-una} a differential comparison at $\qq_{\geq n}$
must record disagreement as an expected consequence of the semantics, with the
unique-name assumption named as its cause. Recording it as a defect would
motivate a repair that could only be effected by falsifying one of the two
semantics.
\end{corollary}

\begin{remark}[Inheritance is not vacuity]
\label{rem:inheritance-guard}
\Cref{prop:inherited-una} makes \textsf{UNIQUE\_NAMES} free, but not
unconditionally safe. If the term-construction map is not injective---if two
distinct individuals are mapped to the same ground term because both carry the
same identifier under the chosen naming scheme---then the inherited assumption
silently identifies them, and the lower bound of
\Cref{prop:counting-split}(i) is computed over a collapsed set. A guard
checking injectivity of the term map on the individuals occurring in the
question is therefore a precondition for the inheritance to be sound.
\end{remark}

\subsection{Bounding a question can make it cheaper}

\begin{proposition}[Bounded paths are non-recursive]
\label{prop:bounded-cheaper}
Let $\qq^\infty$ ask for reachability under a role $r$ and $\qq^{\leq k}$ ask
for reachability in at most $k$ steps. A lowering of $\qq^{\leq k}$ by
unrolling produces $O(k)$ non-recursive clauses, each depending only on the
clause one level below, so evaluation requires no fixpoint. A lowering of
$\qq^\infty$ requires recursion and, on a cyclic corpus, a fixpoint
computation.
\end{proposition}

\begin{proof}
Unrolling defines $r_1 = r$ and
$r_{i+1}(x,z) \leftarrow r_i(x,y), r(y,z)$ for $i < k$, with the goal the union
of the $r_i$. The dependency graph on the $r_i$ is a path and hence acyclic,
so stratified evaluation in $k$ passes suffices. The unbounded question defines
$r^*(x,z) \leftarrow r^*(x,y), r(y,z)$, whose dependency graph has a self-loop.
\end{proof}

\begin{corollary}[Inversion]
\label{cor:inversion}
There exist corpora and engines for which $\qq^{\leq k}$ is answered with
certification while $\qq^\infty$ exceeds the search budget on the same input,
and simultaneously a description logic answers $\qq^\infty$ readily---by a
transitive super-role, per \Cref{rem:no-base-transitive}---while being unable
to \emph{state} $\qq^{\leq k}$ at all, since a transitive role closes
completely and admits no step bound. One word in the question determines which
engine can answer it.
\end{corollary}

\begin{remark}[What $k$ inflates]
\label{rem:k-inflates}
$k$ multiplies the number of clauses, not the size of the search. The
unrolled program grows linearly in $k$ while each clause is evaluated over a
join whose arity is fixed by the clause shape, so the cost is polynomial in
the corpus with the exponent set by that arity rather than by $k$.
\end{remark}

\begin{proposition}[Recursion shape is a cost, not a capability boundary]
\label{prop:recursion-shape}
The right-recursive and left-recursive formulations of transitive closure have
the same least model and hence the same answers, but their evaluation costs
under a goal-directed strategy differ by a factor that grows with the corpus.
\end{proposition}

\begin{proof}
Both programs define the same relation, so their least models coincide by
uniqueness of the least fixpoint of a monotone operator. Under goal-directed
evaluation with a fixed selection rule, the left-recursive form re-derives
prefixes of every partial path before extending it, while the right-recursive
form extends a single frontier; the number of derivation steps therefore
differs, with the ratio increasing in the number of paths.
\end{proof}

\begin{remark}[Assert the shape, not the value]
\label{rem:assert-shape}
Because the two formulations agree on answers, a test asserting the answer set
cannot distinguish them, and a subsequent simplification that replaces one
with the other passes every such test while reintroducing the cost. The test
must assert the emitted \emph{shape}.
\end{remark}

\section{One derivation}
\label{sec:single}

\begin{definition}[View]
\label{def:view}
A \emph{view} is a function from a record set to a presentation. Two views
$v_1, v_2$ \emph{disagree} if there is a record set $S$ and a proposition $p$
about the system such that $v_1(S)$ asserts $p$ and $v_2(S)$ asserts
$\neg p$.
\end{definition}

\begin{theorem}[Single derivation]
\label{thm:single}
Let $d$ be a derivation function from record sets to derived facts, and let
every view factor as $v_i = \pi_i \circ d$ with $\pi_i$ a pure presentation.
Then no two views disagree on any $p$ in the image of $d$.
\end{theorem}

\begin{proof}
Suppose $v_1(S)$ asserts $p$ and $v_2(S)$ asserts $\neg p$ with $p$ derived.
Then $\pi_1(d(S))$ and $\pi_2(d(S))$ disagree about a component of the single
value $d(S)$. Since each $\pi_i$ presents that component without recomputing
it, both read the same value, so they assert the same proposition---
contradiction.
\end{proof}

\begin{corollary}[Duplication admits contradiction]
\label{cor:duplication}
If two views each compute $d$ independently, as $v_i = \pi_i \circ d_i$ with
$d_1 \neq d_2$, then \Cref{thm:single} does not apply and disagreement is
possible with no schema violation, since both outputs are individually well
formed.
\end{corollary}

\Cref{cor:duplication} applies to documents as much as to programs. A figure
environment duplicated between a manuscript and a caption file is two
derivations of one fact, and they drift silently.

\begin{proposition}[A status field is a claim with no clock]
\label{prop:status-clock}
Let a record carry a status $s$ intended to summarise progress, and let $s$ be
set by hand. Then $s$ is a claim about the record's state at the time of
writing, and no subsequent event updates it. In particular a status
\textsf{planned} on an entity for which completed experiments exist is a false
claim that no check over shape can detect.
\end{proposition}

\begin{proof}
Falsity: \textsf{planned} asserts that no run has occurred, and the existence
of a completed experiment referencing the entity refutes it. Undetectability by
shape: the field is of the correct type and drawn from the permitted
enumeration, so schema validation passes. The contradiction is between two
records and is therefore a content property.
\end{proof}

\begin{corollary}[Derive status, or check it against evidence]
\label{cor:derive-status}
Either $s$ is computed by $d$ from the experiment records, in which case
\Cref{thm:single} applies and it cannot contradict them, or $s$ is authored
and a content check must compare it against them. There is no third option in
which an authored status remains true by discipline: nobody sets a status
backwards to a lie, so the errors are all in the direction of understating
progress, and understatement is what makes a coverage table report ``nothing
has been run'' with completed runs on file.
\end{corollary}

\begin{remark}[A gate whose name exceeds its check]
\label{rem:gate-name}
A validation gate named for a strong property while testing a weak consequence
of it is worse than no gate, because its passing is cited as evidence for the
strong property. If a gate is to certify that a derivation used a particular
axiom, it must inspect the axiom's presence; certifying only that the derived
count matched an expectation is satisfied by any run whose arithmetic
coincides, including runs in which the axiom was absent. For the same reason a
count is the weakest of a family of checks: naming the specific entities
expected, and specifically naming a \emph{control} entity that must not appear,
discriminates where a count cannot.
\end{remark}

\begin{remark}[Hand-copied expectations]
\label{rem:hand-copied}
An expectation transcribed by hand into a gate is visible in review; an
expectation parsed from the artefact under test is only as reliable as the
parser, and shares any defect the artefact has. The first is preferable
despite being less automatic.
\end{remark}

\section{Computational validation}
\label{sec:validation}

\subsection{Method}

We implement the constructions of \Cref{sec:setting}--\Cref{sec:single} once,
in a shared module, and state seventeen propositions as executable
experiments. Each experiment reports a verdict in the sense of
\Cref{def:verdict} and a \textsc{pass}/\textsc{fail} against a
\emph{pre-registered} expectation recorded in the experiment's own source
before the run. Where a claim is definitional rather than empirical the
experiment says so in its \texttt{interpretation} field, and where an
experiment's expected outcome is a refusal, the refusal is the \textsc{pass}
condition.

Every experiment that could be satisfied trivially carries a control that is
able to fail. We state this requirement precisely because it is the one most
easily satisfied in appearance only: a control is a control if and only if it
would produce a different result under the negation of the hypothesis, which
is exactly \Cref{cor:control-differs}.

The suite runs on a synthetic knowledge base whose ground truth is known by
construction, together with two engines: a description-logic-like evaluator
implementing open-world certain-answer semantics without a unique-name
assumption, and a rule evaluator implementing least-model semantics over a
Herbrand universe. Both are implemented in the shared module so that the
experiments test the definitions rather than a third-party implementation's
conformance to them. This is a deliberate limitation and we return to it in
\Cref{sec:discussion}.

\subsection{The propositions under test}

\begin{table}[H]
\centering
\small
\begin{tabular}{@{}llp{7.4cm}@{}}
\toprule
\# & Claim & What the experiment does \\
\midrule
01 & \Cref{thm:independence} & Constructs the three witnesses and checks each satisfies two predicates and fails the third. \\
02 & \Cref{cor:conflation} & Shows five distinct situations produce identical empty answer sets under single-output reporting. \\
03 & \Cref{thm:nondegeneracy} & Checks no non-\vans{} verdict carries an answer set, over all questions in the suite. \\
04 & \Cref{prop:counting-split} & Separates $\geq n$, $\leq n$, $=n$ by which of the two requirements each needs. \\
05 & \Cref{rem:una-scope} & Contrasts the sound scoped check against the tempting global-boolean check. \\
06 & \Cref{prop:compare-refusals} & Scores honest refusal against confident wrong zero under a naive and a verdict-aware rule. \\
07 & \Cref{thm:containment} & Checks the containment test fires exactly when it should, over all question--engine pairs. \\
08 & \Cref{prop:under} & Attempts the counting lowering, records its rejection, and records the empty certified answer set obtained by removing the gate. \\
09 & \Cref{prop:half-true} & Shows the conflated alphabet admits no honest capability set, and the split one does. \\
10 & \Cref{prop:nameable} & Runs a differential comparison with a feature absent from the alphabet and confirms the discrepancy is misattributed. \\
11 & \Cref{thm:irredundance}(ii) & Constructs a mis-rooted negation: syntactically distinct programs, identical answers. \\
12 & \Cref{thm:irredundance}(i) & Constructs a dropped constraint: identical programs, and confirms the syntactic check reports the stronger fact. \\
13 & \Cref{prop:agreement} & Confirms two engines agreeing under a vacuous constraint yields no evidence of faithfulness. \\
14 & \Cref{prop:disjoint-detector}, \Cref{cor:localisation} & Injects a genuine violation; compares global inconsistency against a localising fact. \\
15 & \Cref{prop:inherited-una}, \Cref{cor:disagree-result} & Measures the legitimate disagreement at $\qq_{\geq n}$ and attributes it. \\
16 & \Cref{prop:bounded-cheaper}, \Cref{cor:inversion} & Measures evaluation cost of bounded versus unbounded reachability on a cyclic corpus, both engines. \\
17 & \Cref{thm:single}, \Cref{prop:status-clock} & Builds two views over one derivation and two over duplicated derivations; injects a stale status and checks which regime detects it. \\
\bottomrule
\end{tabular}
\caption{The seventeen validation experiments and the claims they test.
Experiments 08, 11, 12 and 15 are ones whose \textsc{pass} condition is a
refusal or a disagreement rather than an answer.}
\label{tab:experiments}
\end{table}

\subsection{Results}

Results are written as one JSON document per experiment together with a
summary document, in \texttt{ckg/validation-incompleteness/results/}. Each
document records the pre-registered expectation, the measured outcome, the
verdict, the control's outcome, and an interpretation stating what the result
does \emph{not} establish. We draw attention to four outcomes.

\paragraph{Experiment 06.} Under a scoring rule counting answers returned, the
engine giving a confident wrong zero outscores the honest refuser; under a
verdict-aware rule the ordering reverses. The magnitude of the reversal is the
ground-truth cardinality, as \Cref{prop:compare-refusals} predicts.

\paragraph{Experiment 08.} The counting lowering is rejected by the fragment
gate. With the gate removed the run terminates and returns an empty answer set
with certification asserted---the answer set of a program whose aggregate never
grounded. This is the configuration \Cref{cor:certified-empty} identifies as
maximally misleading, and it is reachable by removing a check that an audit
conducted against the target language would recommend removing.

\paragraph{Experiment 11.} The mis-rooted negation produces programs that
differ syntactically and answer identically. The syntactic check passes and the
semantic check fires, confirming \Cref{thm:irredundance}(ii): a syntactic guard
cannot catch a semantic vacuity.

\paragraph{Experiment 16.} Bounded reachability is answered with certification
where unbounded reachability on the same cyclic corpus exceeds the search
budget, while the open-world engine answers the unbounded question readily and
returns \vsurf{} on the bounded one. Both directions of \Cref{cor:inversion} are
realised on a single corpus, at one budget, on twelve assertions. The refusal
is \vsurf{} rather than \vexp{} because the model does carry
\textsf{BOUNDED\_PATH}: the limit is in the engine's compiler, not in the
language, and it is exactly the distinction of \Cref{thm:independence} that
lets the two be told apart. Measured step counts for $k = 1,\dots,5$ are
$1,2,3,4,5$, confirming the cost claim of \Cref{prop:bounded-cheaper} that
unrolling scales with the bound and not with the corpus.

\subsection{What the suite does not establish}

The engines are ours, so conformance of any third-party reasoner to
\Cref{def:derivable} is not tested; the corpus is synthetic, so no claim about
biochemical coverage follows; and the timing in Experiment 16 is a property of
our evaluator's strategy, not of any standard. Experiments 03, 07 and 09 are
definitional---they check that an implementation obeys a definition---and
cannot fail except by implementation error, which is what they are for.

\section{Discussion}
\label{sec:discussion}

\subsection{What the framework buys}

The framework converts a class of silent defects into loud ones. Each of the
five situations in \S1.1 now produces a distinct, structured output, and
\Cref{thm:nondegeneracy} makes the distinctness a checkable property rather
than a convention. The two vacuity checks convert a defect that previously
manifested as a correct-looking answer into a verdict, and
\Cref{thm:irredundance} shows that implementing only one of them leaves a
gap. \Cref{thm:single} converts a class of documentation drift into an
impossibility.

The cost is real. A verdict-returning system is more work to build than one
returning answer sets, its outputs are harder to aggregate---one cannot average
a verdict---and its refusals will be read by some consumers as failures. The
last is not a presentational problem but a scoring one, and
\Cref{prop:compare-refusals} identifies the scoring rules under which it
occurs.

\subsection{What the framework does not do}

It does not make a knowledge graph complete, and it is not intended to. It does
not tell you which of the three blockers is cheapest to remove; the unblocking
relation of \Cref{def:unblocking} names candidates and says nothing about
their cost. It does not detect a lowering that is unfaithful in a way that
changes answers without making a constraint inert---a mis-lowering that yields
\emph{different wrong} answers passes both vacuity checks, and only a
comparison against ground truth or against an independent engine catches it,
with \Cref{prop:agreement} limiting what such a comparison establishes.

Most importantly, the framework does not remove the need for judgement about
which questions to ask. A question set assembled to demonstrate a system's
strengths will produce a verdict distribution favourable to that system, and
\Cref{thm:nondegeneracy} makes the individual verdicts honest without making
the aggregate representative. A fixture built to exhibit a property is the
weakest available evidence that anything else has it.

\subsection{Outlook}

Two extensions seem worth pursuing. The first is a cost model over the
unblocking relation, turning \Cref{def:unblocking} from an enumeration into a
ranking: given a blocked question set, which single change to model, corpus,
engine or budget unblocks the most. The second is the treatment of
\emph{partial} derivability---questions determined on part of their domain and
underdetermined elsewhere---which the present binary $\mathrm{Der}$ predicate
collapses, and which is the common case on real corpora where closure holds for
curated entities and fails for imported ones.

\section*{Reproducibility}

The validation suite is deterministic; every stochastic component seeds
explicitly. Each result document records the expectation before the measured
value, and expectations are transcribed by hand into the experiment sources
rather than parsed from the artefacts under test, for the reason given in
\Cref{rem:hand-copied}. Result documents record \texttt{null} where nothing was
measured, and never a zero.

\bibliographystyle{plainnat}
\bibliography{references}

\end{document}
